A procedure is not a tool, and the question ``does it work'' does not reduce to a
benchmark score. What Section~3 produces is a partition and a bound, not a
measurement, so this section states what evaluating such an output can and cannot
mean before Section~5 applies it.

Four questions are answerable, and they are not equally strong. A fifth is not
answerable here and is named as such.

\subsection{Did the bound hold?}

The partition claims that defects corrupting an artifact the oracle consumes are
invisible to it. The claim is testable one defect at a time: build the system at the
defect, run the oracle, and see whether it fires.

The evidence divides, and the division matters. Most calibration in Section~5 was
carried out before the partition existed and for other oracles, so it does not test
this partition. Of the calibration performed afterwards, three pairs were
\emph{selected using} the partition and are therefore confirmatory rather than
independent: the partition said where to look, and looking there succeeded.

Twenty-nine cases are genuinely held out. One was the original: a defect classified as
invisible by the procedure alone---three source lines aligned, no system built---and
checked afterwards against the fix and the oracle's own predicate. Twenty-eight more were
added under a pre-registration discipline, in three batches, and the discipline is the
part worth stating. The frame was fixed by a mechanical query, the sample drawn by a
content-independent rule, and each case classified from the kernel diff alone with a
falsifiable prediction attached. Classifications were committed \emph{before} any
verification ran, in a separate commit from the results, so the ordering is not a
claim about our conduct---it is visible in the artifact's history.

All twenty-eight predictions held, and the composition matters more than the count.
Twenty-one corrupt nothing the oracle consumes, so they answer the step's second
case rather than testing the bound; eighteen of those lie off the
state-equivalence path altogether. Seven tested the blind bound---the fields at issue
were a precision flag, register identities, a pointer type, and stack-slot
metadata, each read from the system's own capture rather than recomputed---and all
seven were confirmed blind.

The first two batches drew from an unfiltered frame and produced no case touching
the predicate at all, leaving the bound's upper arm untested. A third batch was
drawn from a declared stratum---the same frame restricted to fixes whose
\emph{message} mentions pruning or state equivalence---and the stratification is
defensible for a measured reason: in an earlier population, four fixes in forty
carried titles naming a component the diff did not touch, so the message does not
determine the classification the source must still make. That batch supplied the
missing case. A fix replacing the scalar arm's equality test with an
identity-mapping check was classified visible, and the oracle implements exactly
that arm---indeed it is the arm on which the oracle's own calibrated positive
fires. One further case is sharper than either arm: a fix to a second comparison
function, correctly classified as touching the predicate, which the oracle
nonetheless cannot see because it does not model that function. Predicate does not
imply visible, and the held-out set now contains an instance rather than an
assertion.

Two further limits belong here rather than in Section~6, because they bound what
these twenty-eight establish, and the first needs a distinction to be read
correctly. Two separate acts in this study are performed from source. Classifying a
defect from its fix, without building anything, is not a shortcut---it is the
procedure, and doing it any other way would not be testing the procedure at all.
Checking afterwards whether the field that defect touches enters the oracle's
verdict as a consumed value or an independently recomputed one is a different act,
and there a build is stronger evidence than a reading. For twenty-seven of the
twenty-eight, no system was built. For the visible case we built both kernels,
because a reading-only confirmation of the one prediction on the visible side would
have been the weakest link in the study, and what the build returned is more useful
than a confirmation. On the first run the oracle was \emph{silent} on the defective
kernel, although the capture contained exactly the pair the fix's own commit message
describes. The cause was not the predicate: the oracle routes captures to an
era-specific implementation of the comparison by their print format, the 2023 format
is shared with a 2021 era whose comparison legitimately had no identity check, and
the oracle had no arm for the year in between. The source-level reading ``the arm is
implemented'' had been true of one path through the oracle and false of the path this
input took. With the era stated explicitly the oracle fired on the defective kernel
with the predicted arm and was silent on the fixed one, where the pruning event
itself no longer occurs; the same capture under the old routing stays silent. We
count this as a built calibration pair, not as a held-out confirmation, because the
oracle was changed after the first run; and we count the first run as the study's
sharpest result, because it is the axis observed inside the oracle: what the oracle
did was decided by what it read. None of this changes the count. The blind side of
the bound rests on seven held-out cases; the visible side rests on one, drawn from a
stratified batch. The build hardens that one case; it does not add a second, and the
asymmetry is permanent in this study.

The second limit is a wording error of ours. One prediction's falsification
condition, as written in the first batch, conflated ``modelled in the oracle'' with
``independently recomputed by the oracle''; the substantive criterion is the second,
the wording said the first, and a stricter reading would score that case ambiguous
rather than confirmed. The condition was corrected before the second batch, and both
readings are recorded.

\subsection{Can a reader check it?}

For a procedure whose output is a classification, auditability is a more appropriate
evaluation than accuracy, and a stronger one to offer. Every decision in the
partition cites the fix that anchors it and the location in that fix which settles
the question. The unit of verification is one commit: a reader who doubts a row can
open it and re-derive the decision in minutes, without the system, the corpus, or
any part of our infrastructure.

This is what a classification procedure can offer in place of a score. The claim is
not ``trust the partition'' but ``check any row of it''---and it now holds for every partition
in the paper. It did not always. The first flagship partition, of forty fixes, was
classified with anchors recorded in an engineering log but without retaining the
per-commit list as one artifact; when a reviewer asked whether its rows could be
checked, they could not be enumerated, and the partition was rebuilt from a
mechanical frame with both raters' files kept (Section~5.3). The lesson is procedural
and we state it as one: a partition's value as evidence depends on keeping the list,
and the cost of not keeping it was a headline number that had to be withdrawn.

\subsection{Did it change what was done?}

A bound that changes no decision has no value regardless of its correctness, and
the strongest evidence that this one does is that two oracles in this project were
built \emph{without} it and paid for that. Oracle $B$ was designed as an independent
reimplementation of the state-equivalence predicate---independent in the only sense
then in use---and its evidential entanglement was discovered afterwards, by reading
the first partition, at which point close to half of the population it had been built
to catch turned out to have been unreachable from the day it was written. A store-location channel in
the same project judged runtime writes against the verifier's own printed claim of
where each store would land; nine successive false findings were traced, one by one,
to that consumed claim rather than to any predicate, and the verdict had to be moved
to an input the verifier did not produce. With the axis in hand neither design would
have been chosen as it was. Both are recorded in the engineering log as discoveries,
which is the point: they were not obvious at the time.

A smaller instance is recorded as well. Before the partition, the next calibration target had been
chosen by reading fix summaries. The partition classified that target as
structurally invisible to the oracle in question and redirected the choice to a
different defect. The redirected target produced a positive calibration; the
original would have required building two versions of the system to discover that it
could not have fired. The procedure replaced a build with a reading.

\subsection{Does it generalize?}

It is not shown here, and the reason is itself a result rather than an omission.
Validating the bound outside one system requires a defect population classified by
whether an oracle could have observed each defect. No such data exists: published
taxonomies organize defects by symptom and by configuration; a technique's own bug
list measures its sighted side and cannot bound anyone's blind one---including its
own; and the scope statements techniques do publish are component-scoped rather than
input-scoped, which Section~7 shows is not the same partition. Section~7 names the study that would supply the missing data and the design
it must have.

\subsection{Reproducibility, measured once}

Two analysts applying Section~3 to the same population should produce the same
partition. We measured this on the twenty-eight held-out cases with a second rater
that had no access to the first rater's labels: an independent model instance given
only the question of Section~3.3, a factual description of the oracle's consumption
set, and the diffs. It is not a human rater, and what the experiment measures is
correspondingly narrower---whether an independent reader of the same source, asked
the same question, gives the same answer.

On the held-out cases, raw agreement was 25 of 28 and Cohen's $\kappa = 0.71$,
substantial but not near-perfect. On the rebuilt flagship partition of fifty-seven,
agreement was 53 of 57 and $\kappa = 0.86$. Both are above the threshold we
pre-registered as falsifying the auditability claim. The second measurement carries a
caveat the first does not: the first rater read the fifty-seven diffs as truncated
excerpts and the second read them in full, and two of the four disagreements are
plainly the first rater's under-reading. Two details matter more than the coefficient. On the held-out cases the second rater
answered \emph{Undecided} twice; the first rater answered it zero times in the
sixty-eight classifications up to that point, which Section~6 had already flagged as
suspicious, and the flag now has a measurement behind it. On the rebuilt flagship
partition each rater answered it once, on different cases. And the three disagreements are not
scattered: all three are fixes that write a \emph{stack slot}, which this oracle does
not consume at all. The first rater called them blind-by-corruption (Yes); the second
called them not-consumed (No) or Undecided. Both readers agree the oracle is blind to
them; they disagree on which of the step's answers that blindness belongs to, and the
step's wording is what is at fault: its ``No'' covers both a defect the oracle could
in principle observe and a defect in an artifact the oracle never reads, and only the
first is potentially visible. Section~3.3 is corrected accordingly. This is also the
right way to read what the model rater delivered and what it did not. The coefficient
does not rescue auditability; it makes auditability measurable. What the experiment
actually produced is the seam, and that finding does not depend on the rater being
human, because the defect it exposed is in the step's wording, not in any reader of
it. A human study would supply an effect size for the corrected step; it would not
supply the correction, which is already in hand.
